\documentclass[12pt]{article}
%^ Start of document
\usepackage{biblatex} %Imports biblatex package
\usepackage{graphicx} % Required for inserting images
\usepackage{authblk}
\usepackage[colorlinks=true, urlcolor=blue, linkcolor=red]{hyperref}
\usepackage{subfig}
\usepackage{float}
\usepackage{tabularx}
\usepackage{multirow}
\usepackage[paperheight=11in,
   paperwidth=8.5in,
   top=20mm,
   bottom=20mm,
   left=40mm,
   right=40mm]{geometry}
\usepackage{moresize}
\usepackage{tikz}
\usepackage{xcolor}
\usepackage{physics}
\usepackage{amssymb}
\usepackage{mathrsfs}
\usepackage{amsmath}
\usepackage{ragged2e}
\usepackage{comment}
\usepackage{xcolor}
\addbibresource{Electro.bib}
\begin{document}

\title{Classical Electrodynamics Exam 1}

\author[1]{Zachary Tullier}
\affil[1]{Physics Masters Student\\University of Houston - Clear Lake\\Physics Department\\Houston, Texas\\U.S}

\maketitle

\section{Laplace Equation in Spherical Coordinates} \label{problem1}
    In this problem we will consider a sphere with radial distance ($r$), polar angle that the point makes with the z-axis ($\theta$), and the azimuthal angle in the x-y plan relative to the x-axis ($\phi$). 
    \newline The Laplace equation is defined as
    \begin{equation} \label{eq: 1.1}
        \nabla^2 \Phi = 0
    \end{equation}
    Expand the $\nabla$ and convert to spherical coordinates
    % \newline \textcolor{red}{Do in between}
    \begin{equation} \label{eq: 1.2}
        \frac{1}{r} \frac{\partial^2}{\partial r^2} (r \Phi) + \frac{1}{r^2 sin(\theta)} \frac{\partial}{\partial \theta} \left(sin(\theta) \frac{\partial \Phi}{\partial \theta} \right) + \frac{1}{r^2 sin^2(\theta)} \frac{\partial^2 \Phi}{\partial \phi^2} = 0
    \end{equation}
    Use method of separation of variables by trying a solution of the form
    \begin{equation} \label{eq: 1.3}
        \Phi(r,\theta,\phi) = \frac{R(r)}{r} P(\theta) Q(\phi)
    \end{equation}
    Substitute equation \ref{eq: 1.3} into equation \ref{eq: 1.2}
    \begin{equation} \label{eq: 1.4}
        \begin{split}
        \frac{1}{r} \frac{\partial^2}{\partial r^2} \left(r \frac{R(r)}{r} P(\theta) Q(\phi) \right) 
        \\ + \frac{1}{r^2 sin(\theta)} \frac{\partial}{\partial \theta} \left(sin(\theta) \frac{\partial \frac{R(r)}{r} P(\theta) Q(\phi)}{\partial \theta}\right) 
        \\+ \frac{1}{r^2 sin^2(\theta)} \frac{\partial^2 \frac{R(r)}{r} P(\theta) Q(\phi)}{\partial \phi^2} = 0
        \end{split}
    \end{equation}
    We want to rearrange the equation to match a differential equation; therefore, we multiply both sides by $\frac{r^3 sin^2(\theta)}{R(r) P(\theta) Q(\phi)}$
    % \newline \textcolor{red}{Do in between}
    \begin{equation} \label{eq: 1.5}
        \frac{r^2 sin^2(\theta)}{R(r)} \frac{d^2 R(r)}{dr^2} + \frac{sin(\theta)}{P(\theta)} \frac{d}{d\theta}\left(sin(\theta) \frac{dP(\theta)}{d\theta}\right) + \frac{1}{Q(\phi)} \frac{d^2Q(\phi)}{d\phi^2} = 0
    \end{equation}
    The last term is solely dependent on $\phi$ so it can be equal to a constant ($m$)
    \begin{equation} \label{eq: 1.6}
        \frac{r^2 sin^2(\theta)}{R(r)} \frac{d^2 R(r)}{dr^2} + \frac{sin(\theta)}{P(\theta)} \frac{d}{d\theta}\left(sin(\theta) \frac{dP(\theta)}{d\theta}\right) - m^2 = 0
    \end{equation}
    Where 
    \begin{equation} \label{eq: 1.7}
         -m^2 = \frac{1}{Q(\phi)} \frac{d^2Q(\phi)}{d\phi^2}
    \end{equation}
    Rearrange to get common solution
    % \newline \textcolor{red}{How do I get common solution}
    \begin{equation} \label{eq: 1.8}
        \frac{d^2Q(\phi)}{d\phi^2} = -m^2 Q(\phi)
    \end{equation}
    If $m\neq0$ the common solution is 
    \begin{equation} \label{eq: 1.9}
        Q(\phi) = A_m e^{im\phi} + B_m e^{-im\phi}
    \end{equation}
    If $m=0$ the common solution is 
    \begin{equation} \label{eq: 1.10}
        Q(\phi) = A_m + B_m \phi
    \end{equation}
    Now we focus on the Original equation \ref{eq: 1.6}
    \newline  We want to rearrange the equation to match a differential equation; therefore, we multiply both sides by $sin^2(\theta)$ 
    \begin{equation} \label{eq: 1.11}
        \frac{r^2}{R(r)} \frac{d^2 R(r)}{dr^2} + \frac{1}{sin(\theta) P(\theta)} \frac{d}{d\theta}\left(sin(\theta) \frac{dP(\theta)}{d\theta}\right) - \frac{m^2}{sin^2\theta)} = 0
    \end{equation}
    The first and last terms are independent so they can be set to constants
    \begin{equation} \label{eq: 1.12}
        \frac{r^2}{R(r)} \frac{d^2 R(r)}{dr^2} - l (l+1) = 0
    \end{equation}
    Where 
    \begin{equation} \label{eq: 1.13}
        -l(l+1) = \frac{1}{sin(\theta) P(\theta)} \frac{d}{d\theta}\left(sin(\theta) \frac{dP(\theta)}{d\theta}\right) - \frac{m^2}{sin^2\theta)}
    \end{equation}
    Put equation \ref{eq: 1.12} into more intuitive forms
    \begin{equation} \label{eq: 1.14}
        \frac{d^2 R(r)}{dr^2} = l(l+1) \frac{R(r)}{r^2} 
    \end{equation}
    Put equation \ref{eq: 1.13} into more intuitive forms
    \begin{equation} \label{eq: 1.15}
        \frac{d}{d\theta} \left(sin(\theta) \frac{d P(\theta)}{d\theta}\right) + \left(l(l+1)-\frac{m^2}{sin^2(\theta)}\right) sin(\theta) P(\theta) = 0
    \end{equation}
    Reduce equation \ref{eq: 1.15} to simpler form by substituting $x=cos(\theta)$ and $\frac{d}{d\theta} = -\sqrt{1-x} \frac{d}{dx}$
    \begin{equation} \label{eq: 1.16}
        \frac{d}{dx} \left((1-x^2) \frac{d P(X)}{dx}\right) + \left(l(l+1)-\frac{m^2}{1-x^2}\right) P(x) = 0
    \end{equation}
    Equation \ref{eq: 1.14} by substituting $R(r) = r^{\alpha}$
    \begin{equation} \label{eq: 1.17}
        \alpha (\alpha - 1) r^{\alpha - 2} = (l+1) r^{\alpha - 2}
    \end{equation}    
    % \textcolor{red}{Do in between}
    \begin{equation} \label{eq: 1.18}
        \alpha^2 - \alpha - 1) - l(l+1) = 0
    \end{equation}        
    % \textcolor{red}{Do in between}
    Solve Polynomial
    \begin{equation} \label{eq: 1.19}
        \begin{split}
            \alpha = l + 1
            \\ \alpha = -l
        \end{split}        
    \end{equation}
    The General solution is 
    \begin{equation} \label{eq: 1.20}
        R(r) = A_l r^{l+1} + B_l r^{-1}
    \end{equation} 

\section{Ordinary Legendre Polynomials} \label{problem2}
    The general solution for the laplace equation in spherical coordinates can be formulated by substituting equations \ref{eq: 1.9}, \ref{eq: 1.10}, and \ref{eq: 1.20} into equation \ref{eq: 1.3}
    \begin{equation} \label{eq: 2.1}
        \Phi(r,\theta, \phi) = \sum_m \sum_l (A_l r^l + B_l r^{-l-1}) (A_m d^{im \phi} + B_m e^{-im \phi}) P_l^m(cos(\theta))
    \end{equation}     
    We will assume $m=0$ and use equation \ref{eq: 1.14}
    \begin{equation} \label{eq: 2.2}
        \frac{d}{dx} \left((1-x^2) \frac{d P(X)}{dx}\right) + l(l+1) P(x) = 0
    \end{equation}
    Guess the power series solution 
    % \newline \textcolor{red}{How do we get this "guess}
    \begin{equation} \label{eq: 2.3}
        P(x) = \sum_{j=0}^{\infty} a_j x^{j+\alpha}
    \end{equation}    
    % \textcolor{red}{Do Inbetween}
    \begin{equation} \label{eq: 2.4}
        \sum_{j=0}^{\infty} (j+\alpha) (j + \alpha - 1) a_j x^{j+\alpha-2} - (j+\alpha) (j + \alpha + 1) a_j x^{j+\alpha} + l(l+1) a_j x^{j+\alpha} = 0
    \end{equation}        
    Consolidate
    \begin{equation} \label{eq: 2.5}
        \sum_{j=0}^{\infty} (j+\alpha) (j + \alpha - 1) a_j x^{j+\alpha-2} - \sum_{j=0}^{\infty} ((j+\alpha) (j + \alpha + 1) + l(l+1)) a_j x^{j+\alpha} = 0
    \end{equation}            
    Remove first two powers of x then combine
    % \newline \textcolor{red}{Do Inbetween}
    \begin{equation} \label{eq: 2.6}
        \begin{split}
        (\alpha) (\alpha - 1) a_0 x^{\alpha -2)} + (1 + \alpha) (\alpha) a_1 x^{\alpha - 2} 
        \\ - \sum_{j=0}^{\infty} ((j + 2 + \alpha) (j + \alpha + 1) a_{j + 2} - ((j + \alpha) (j + \alpha + 1) 
        \\- ((j + \alpha) (j + \alpha + 1) a_j) x^{j+ \alpha}
        \end{split}
    \end{equation}            
    This must hold for all values of x so
    \begin{equation} \label{eq: 2.7}
        (\alpha) (\alpha - 1) a_0 = 0
    \end{equation}            
    \begin{equation} \label{eq: 2.8}
        (1 + \alpha) (\alpha) a_1 = 0
    \end{equation}            
    Substitute 
    % \newline \textcolor{red}{Do Inbetween}
    \begin{equation} \label{eq: 2.9}
        (j + \alpha + 2) (j + \alpha + 1) a_{j + 2} - ((j + \alpha) (j + \alpha + 1 - l(l+1))) = 0
    \end{equation}  
    The third equation is solved to yield the recurrence relation which gives all the remaining terms from the last one. Rerrange equation
    \begin{equation} \label{eq: 2.10}
        a_{j+2} = \frac{(j+\alpha)(j+\alpha+1) - l(l+1)}{(j + \alpha + 2) (j + \alpha + 1} a_j
    \end{equation}      
    $a_1 = 0$; therefore, $a_{odd} = 0$
    \begin{equation} \label{eq: 2.11}
        P_l(x) = \sum_{j=0,even}^{\infty} a_j x^{j+\alpha}
    \end{equation}  
    Converges at $x=1$
    \begin{equation} \label{eq: 2.12}
        \frac{(j_{max}+\alpha)(j_{max}+\alpha+1) - l(l+1)}{(j_{max} + \alpha + 2) (j_{max} + \alpha + 1} =0
    \end{equation}      
    If $a=0$
    \begin{equation} \label{eq: 2.13}
        j_{max} (j_{max} + 1) - l(l+1) =0
    \end{equation}      
    Therefore 
    \begin{equation} \label{eq: 2.14}
        j_{max}=l
    \end{equation}          
    Substitute 
    \begin{equation} \label{eq: 2.15}
        P_l(x) = \sum_{j=0,even}^{l} a_j x^{j}
    \end{equation}      
    Therefore 
    \begin{equation} \label{eq: 2.16}
        P_l(x) = \sum_{j=0}^{\frac{l}{2}} a_{2j }x^{2j}
    \end{equation}    
    Rewrite sums over all integers by doubling indices. If l is even
    \begin{equation} \label{eq: 2.17}
        a_{j+2} = j \frac{(j+1) - l(l+1)}{(j + 2) (j + 1} a_j
    \end{equation} 
    If $\alpha = 1$
    \begin{equation} \label{eq: 2.18}
        (j_{max} + 2)(j_{max} + 1) - l(l+1) =0
    \end{equation} 
    Therefore
    \begin{equation} \label{eq: 2.19}
        j_{max}=l-1
    \end{equation} 
    Substitute 
    \begin{equation} \label{eq: 2.20}
        P_l(x) = \sum_{j=0,even}^{l-1} a_j x^{j+1}
    \end{equation} 
    Rewrite sums over all integers by doubling indices. If l is odd
    \begin{equation} \label{eq: 2.21}
        a_{j+2} = \frac{(j+1) (j+2) - l(l+1)}{(j + 3) (j + 2)} a_j
    \end{equation}     

\section{Laplace Equation Solution for Azimuthal Symmetry} \label{problem3}
    Azimuthal symmetry is when $m=0$, and the region of valid solution includes the entire $2\pi$ radian sweep of $\phi$. The solution will depend on azimuth angle.
    \newline In this problem we will be considering a sphere with radius $a$ with a potential ($V(\theta)$ on its surface and we wish to find the potential everywhere inside the sphere. 
    \newline We begin with the general solution to the Laplace Equation in spherical coordinates.
    \begin{equation} \label{eq: 4.1}
        \Phi(r,\theta, \phi) = \sum_{m=0}^{\infty}  \sum_{l=0}^{\infty} (A_l r^l + B_l r^{-l-1}) (A_m  e^{i m \phi} + B_m e^{-i m \phi}) P_l^m (cos(\theta))
    \end{equation}
    For the special case that $m=0$
    \begin{equation} \label{eq: 4.2}
        \Phi(r,\theta, \phi) = \sum_{l=0}^{\infty} (A_l r^l + B_l r^{-l-1}) P_l (cos(\theta))
    \end{equation}
    Because $2\pi$ radian sweep of $\phi$ of a sphere includes the origin, $B_l = 0$ so that the solution does not blow up to $\infty$
    \begin{equation} \label{eq: 4.3}
        \Phi(r,\theta, \phi) = \sum_{l=0}^{\infty} (A_l r^l) P_l (cos(\theta))
    \end{equation}
    At the surface ($r=a$) we dictate the electric potential ($\Phi = V(\theta)$)
    \begin{equation} \label{eq: 4.4}
        V(\theta) = \sum_{l=0}^{\infty} (A_l r^l) P_l (cos(\theta))
    \end{equation}
    We need to manipulate this equation so that we can use the orthogonality condition. Multiply by $P_{l'}(cos(\theta)) sin(\theta)$ and integrate. The integration is over $0 \xrightarrow{} \pi$ because the orthogonality condition requires it. 
    \begin{equation} \label{eq: 4.5}
        \int_0^{\pi} V(\theta) P_{l'}(cos(\theta)) sin(\theta) d\theta = \sum_{l=0}^\infty A_l a^l \int_0^{\pi} P_{l'}(cos(\theta)) P_l(cos(\theta)) sin(\theta) d\theta
    \end{equation}
    Using the Orthogonality condition
    \begin{equation} \label{eq: 4.6}
        \int_{-1}^{1} P_{l'}(x) P_l(x) dx = \frac{2 \delta_{l'l}}{2l + 1}
    \end{equation}
    Where $x=cos(\theta)$
    \begin{equation} \label{eq: 4.7}
        \int_0^{\pi} V(\theta) P_{l'}(cos(\theta)) sin(\theta) d\theta = \sum_{l=0}^\infty A_l a^l \frac{2 \delta_{l'l}}{2l + 1}
    \end{equation}
    Isolate $A_l$. The $\sum$ will be considered at a higher level equation; therefore, it can be removed
    \begin{equation} \label{eq: 4.8}
        A_l = \frac{2l + 1}{2 a^l \delta_{l'l}}\int_0^{\pi} V(\theta) P_{l'}(cos(\theta)) sin(\theta) d\theta
    \end{equation}
    $l$ and $l'$ are the equivalent: therefore, $\delta_{l'l} = 1$
    \begin{equation} \label{eq: 4.9}
        \boxed{\boldsymbol{A_l = \frac{2l + 1}{2 a^l }\int_0^{\pi} V(\theta) P_l(cos(\theta)) sin(\theta) d\theta}}
    \end{equation}
    \begin{gather*}
        \boxed{\boldsymbol{\Phi(r,\theta, \phi) = \sum_{l=0}^{\infty} (A_l r^l) P_l (cos(\theta))}}
    \end{gather*}
    Shown above, Equation \ref{eq: 4.3} with equation \ref{eq: 4.9} are the general solution for this problem. 

\section{Unit point charge potential expansion in Legendre Polynomials} \label{problem4}
    In this problem we will be considering unit point charge on the $z$ axis with a Azimuthal symmetrical charged sphere with radius $a$ and $m=0$.
    We will begin with the generic form of a unit point charge at the origin
    % \newline \textcolor{red}{Get the background of the equation from lec3 section 2}
    \begin{equation} \label{eq: 5.1}
        \Phi(r,\theta) = \frac{1}{\abs{\vec{r} - \vec{r_0}}} = \frac{1}{\sqrt{r^2 + r_0^2 - 2 r r_0 cos(\theta)}}
    \end{equation} 
    We will also begin with the general solution to the laplace equation in spherical coordinates
    \begin{equation} \label{5.2}
        \Phi(r,\theta, \phi) = \sum_{m=0}^{\infty}  \sum_{l=0}^{\infty} (A_l r^l + B_l r^{-l-1}) (A_m  e^{i m \phi} + B_m e^{-i m \phi} P_l^m (cos(\theta))
    \end{equation}
    For the special case that $m=0$
        \begin{equation} \label{eq: 5.3}
        \Phi(r,\theta, \phi) = \sum_{l=0}^{\infty} (A_l r^l + B_l r^{-l-1}) P_l (cos(\theta))
    \end{equation}
    Combine equation \ref{eq: 5.1} and equation \ref{eq: 5.3}
    \begin{equation} \label{eq: 5.4}
        \frac{1}{\abs{\vec{r} - \vec{r_0}}} = \sum_{l=0}^{\infty} (A_l r^l + B_l r^{-l-1}) P_l (cos(\theta))
    \end{equation}
    This equation must hold for all $\theta$; therefore, to find the coefficients $A_l$ and $B_l$, we dictate that $\theta=0$
    \begin{equation} \label{eq: 5.5}
        \frac{1}{\abs{\vec{r} - \vec{r_0}}} = \sum_{l=0}^{\infty} (A_l r^l + B_l r^{-l-1})
    \end{equation}
    % \textcolor{red}{What is $P_l(1)$}
    Because the general mathematical formulation has two components that converge in different directions we will split this into two regions, $r<r_0$ and $r>r_0$.
    \newline For $r<r_0$ (the observation point is closer to the origin than the charge)
    \begin{equation} \label{eq: 5.6}
        \frac{1}{r_0 - r} = \sum_{l=0}^{\infty} (A_l r^l + B_l r^{-l-1})
    \end{equation}
    We want to rearrange the equation so that we can use the geometric series rule, which only works for $s<1$; therefore, multiply both sides by $r_0$. 
    % \textcolor{red}{How does multiplication end up like this?}
    \begin{equation} \label{eq: 5.7}
        \frac{1}{1 - \frac{r}{r_0}} = \sum_{l=0}^{\infty} (A_l r_0 r^l + B_l r_0 r^{-l-1})
    \end{equation}
    Expand the left side using geometric series rule: $\frac{1}{1-s} = \sum_{l=0}^{\infty} s^l$
    \begin{equation} \label{eq: 5.8}
        \sum_{l=0}^{\infty} {\frac{r}{r_0}}^l =\sum_{l=0}^{\infty} (A_l r_0 r^l + B_l r_0 r^{-l-1})
    \end{equation}
    The left side converges to infinity; therefore the right side also has to converge to infinity. $r^{-l-1}$ converges to zero; therefore, $B_l=0$.
    \begin{equation} \label{eq: 5.9}
        {\frac{r}{r_0}}^l = (A_l r_0 r^l)
    \end{equation}
    Isolate $A_l$
    \begin{equation} \label{eq: 5.10}
        A_l = \frac{r^l}{r_0^l r_0 r^l} = \frac{1}{r_0^{l+1}}
    \end{equation}
    Substitute equation \ref{eq: 5.10} into equation \ref{eq: 5.4}. We will still consider $B_l=0$. 
    \newline The final expanded solution for electric potential if $\boxed{r<r_0}$
    \begin{equation} \label{eq: 5.11}
        \boxed{\boldsymbol{\Phi(r,\theta) = \frac{1}{\abs{\vec{r} - \vec{r_0}}} = \sum_{l=0}^{\infty} (\frac{r^l}{r_0^{l+1}}) P_l (cos(\theta))}}
    \end{equation}
    Using the equation \ref{eq: 5.5} For $r>r_0$ (The observation point is further away from the origin than the charge)
    \newline We can move the r's around in the denominator because of the absolute value property
    \begin{equation} \label{eq: 5.12}
        \frac{1}{r - r_0} = \sum_{l=0}^{\infty} (A_l r^l + B_l r^{-l-1})
    \end{equation}
    We want to rearrange the equation so that we can use the geometric series rule, which only works for $s<1$; therefore, multiply both sides by $r$.
    \begin{equation} \label{eq: 5.13}
        \frac{1}{1 - \frac{r_0}{r}} = \sum_{l=0}^{\infty} (A_l r_0 r^l + B_l r_0 r^{-l-1})
    \end{equation}
    Expand the left side using geometric series rule: $\frac{1}{1-s} = \sum_{l=0}^{\infty} s^l$
    \begin{equation} \label{eq: 5.14}
        \sum_{l=0}^{\infty} {\frac{r_0}{r}}^l =\sum_{l=0}^{\infty} (A_l r_0 r^l + B_l r_0 r^{-l-1})
    \end{equation}
    The left side converges to 0; therefore the right side also has to converge to 0. $r^l$ converges to infinity; therefore, $A_l=0$.
    \begin{equation} \label{eq: 5.15}
        {\frac{r}{r_0}}^l = (B_l r_0 r^{-l-1})
    \end{equation}
    Isolate $B_l$
    \begin{equation} \label{eq: 5.16}
        B_l = \frac{r^l r_0^l}{r^l} = r_0^l
    \end{equation}
    Substitute equation \ref{eq: 5.16} into equation \ref{eq: 5.4}. We will still consider $A_l=0$. 
    \newline The final expanded solution for electric potential if $\boxed{r>r_0}$
    \begin{equation} \label{eq: 5.17}
        \boxed{\boldsymbol{\Phi(r,\theta) = \frac{1}{\abs{\vec{r} - \vec{r_0}}} = \sum_{l=0}^{\infty} \left(\frac{r_0^l}{r^{l+1}}\right) P_l (cos(\theta))}}
    \end{equation}

\section{General Leibnitz Rule}
    The equation that describes the nth derivative of $fg$ was named after Gottfriend Wilhelm Leibniz. Leibniz was a German Polymath who is credited with inventing calculus as well as other branches of mathematics. 
    \newline This rule is typically written as
    \begin{equation} \label{eq: 6.1}
        {(fg)}^{(n)} = \sum_{k=0}^n \begin{pmatrix} n\\ k \end{pmatrix} f^{n-k} g^k
    \end{equation}
    where $\begin{pmatrix} n\\ k \end{pmatrix} = \frac{n!}{k!(n-k)}$ is the binomial coefficient 
    \newline and $f^j$ denotes the jth derivative of f
    \newline This statement must hold for $n+1$
    \begin{equation} \label{eq: 6.2}
        {(fg)}^{(n+1)} = \left(\sum_{k=0}^{n} \begin{pmatrix} n\\ k \end{pmatrix} f^{(n-k)} g^{(k)}\right)'
    \end{equation}
    \begin{equation} \label{eq: 6.3}
        {(fg)}^{(n+1)} = \sum_{k=0}^{n} \begin{pmatrix} n\\ k \end{pmatrix} f^{(n+1-k)} g^{(k)} + \sum_{k=0}^{n} \begin{pmatrix} n\\ k \end{pmatrix} f^{(n-k)} g^{(k+1)}
    \end{equation}
    \begin{equation} \label{eq: 6.4}
        \begin{split}
        {(fg)}^{(n+1)} = \sum_{k=0}^{n} \begin{pmatrix} n\\ k \end{pmatrix} f^{(n+1-k)} g^{(k)} 
        \\+ \sum_{k=0}^{n+1} \begin{pmatrix} n\\ k-1 \end{pmatrix} f^{(n+1-k)} g^{(k)}
        \end{split}
    \end{equation}
    \begin{equation} \label{eq: 6.5}
        \begin{split}
        {(fg)}^{(n+1)} = \begin{pmatrix} n\\ 0 \end{pmatrix} f^{(n+1)} g^{(0)} 
        \\+ \sum_{k=1}^{n} \begin{pmatrix} n\\ k \end{pmatrix} f^{(n+1-k)} g^{(k)} 
        \\+ \sum_{k=1}^{n} \begin{pmatrix} n\\ k-1 \end{pmatrix} f^{(n+1-k)} g^{(k)} + \begin{pmatrix} n\\ n \end{pmatrix} f^{(0)} g^{(n+1)}
        \end{split}
    \end{equation}
    \begin{equation} \label{eq: 6.6}
        \begin{split}
        {(fg)}^{(n+1)} = \begin{pmatrix} n+1\\ 0 \end{pmatrix} f^{(n+1)} g^{(0)} + \left(\sum_{k=1}^{n} \left(\begin{pmatrix} n\\ k-1 \end{pmatrix} + \begin{pmatrix} n\\ k \end{pmatrix}\right) f^{(n+1-k)} g^{(k)}\right) 
        \\+ \begin{pmatrix} n+1\\ n+1 \end{pmatrix} f^{(0)} g^{(n+1)}
        \end{split}
    \end{equation}
    \begin{equation} \label{eq: 6.7}
        \begin{split}
        {(fg)}^{(n+1)} = \begin{pmatrix} n+1\\ 0 \end{pmatrix} f^{(n+1)} g^{(0)} + \sum_{k=1}^{n} \begin{pmatrix} n+1\\ k \end{pmatrix} f^{(n+1-k)} g^{(k)}) 
        \\ + \begin{pmatrix} n+1\\ n+1 \end{pmatrix} f^{(0)} g^{(n+1)}
        \end{split}
    \end{equation}
    \begin{equation} \label{eq: 6.8}
        {(fg)}^{(n+1)} = \sum_{k=0}^{n+1} \begin{pmatrix} n+1\\ k \end{pmatrix} f^{(n+1-k)} g^{(k)}
    \end{equation}


\section{Associated Legendre Function}
    We begin with the General solution for the laplace equation in spherical coordinates
    \begin{equation} \label{eq: 7.1}
        \frac{d}{dx} \left((1-x^2) \frac{dP_l^m(x)}{dx}\right) + \left(l(l+1) - \frac{m^2}{1-x^2}\right) P_l^m(x) = 0
    \end{equation}
    Where $x=cos(\theta)$
    \newline Guess the solution to be 
    \begin{equation} \label{eq: 7.2}
        P_l^m(x) = {(1-x^2)}^{\frac{m}{2}} \frac{d^m}{dx^m} g(x)
    \end{equation}
    Substitute equation \ref{eq: 7.2} into equation \ref{eq: 7.1}
    \begin{equation} \label{eq: 7.3}
        \begin{split}
        \frac{d}{dx} \left((1-x^2) \frac{d\left({(1-x^2)}^{\frac{m}{2}} \frac{d^m}{dx^m} g(x)\right)}{dx}\right) 
        \\+ \left(l(l+1) - \frac{m^2}{1-x^2}\right) {(1-x^2)}^{\frac{m}{2}} \frac{d^m}{dx^m} g(x) = 0
        \end{split}
    \end{equation}
    Partially First derivatives
    \begin{equation} \label{eq: 7.4}
        \begin{split}        
        \frac{d}{dx} \left((1-x^2) \left(-mx{(1-x^2)}^{\frac{m}{2}-1} \frac{d^{m+1}}{dx^{m+1}} g(x)\right)\right) 
        \\+ \left(l(l+1) - \frac{m^2}{1-x^2}\right) {(1-x^2)}^{\frac{m}{2}} \frac{d^m}{dx^m} g(x) = 0
        \end{split}
    \end{equation}
    Use product rule to evaluate derivatives
    \begin{equation} \label{eq: 7.5}
        \begin{split}
        \left(-mx{(1-x^2)}^{\frac{m}{2}} \frac{d^{m+1}}{dx^{m+1}} g(x)\right) + m^2 \frac{d^m}{dx^m}g(x) x^2 {(1-x^2)}^{\frac{m}{2}-1} - m \frac{d^m}{dx^m} g(x) {(1-x^2)}^{\frac{m}{2}} 
        \\+ {(1-x^2)}^{\frac{m}{2}+1} (\frac{d^{m+2}}{dx^{m+2}} - (m+2) x {(1-x^2)}^{\frac{m}{2}} \frac{d^{m+1}}{dx^{m+1}} g(x) 
        \\+ l(l+1){(1-x^2)}^{\frac{m}{2}} \frac{d^m}{dx^m} g(x) - m^2 {(1-x^2)}^{\frac{m}{2}-1} \frac{d^m}{dx^m} g(x) = 0
        \end{split}
    \end{equation}
    Simplify
    \begin{equation} \label{eq: 7.6}
        (1-x^2) (\frac{d^m}{dx^m} \frac{d^2}{dx^2} g(x)) - 2(m+1)x \frac{d^{m+1}}{dx^{m+1}} g(x) + (l(l+1) - m(m+1)) \frac{d^m}{dx^m} g(x)
    \end{equation}
    % \newline \textcolor{red}{Not sure what else to do here. This is a proof of the equation, but it isn't how the equation came to be. And I can't figure out how they use the Leibnitz Rule}

\section{Spherical Harmonics}
    We begin with the full legendre equation
    \begin{equation} \label{eq: 8.1}
        \begin{split}
        \Phi(r,\theta, \phi) = \sum_{m=0,l} (A_l r^l + B_l r^{-l-1}) (A_{m0} + B_{m0}\phi) P_l^{m=0}(cos(\theta)) 
        \\+ \sum_{m\neq0, l} (A_l r^l + B_l r^{-l-1}) (A_m e^{i m \phi} + B_m e^{-i m \phi}) P_l^m (cos(\theta))
        \end{split}
    \end{equation}    
    For this case we will fix the angle, $\phi = $Constant; therefore, $A_{m0} = B_{m0} = 0$
    \newline $m=0$
    % \newline \textcolor{red}{Not sure why}
    \begin{equation} \label{eq: 8.2}
        \Phi(r,\theta, \phi) = \sum_{m\neq0, l} (A_l r^l + B_l r^{-l-1}) (A_m e^{i m \phi} + B_m e^{-i m \phi}) P_l^m (cos(\theta))
    \end{equation}        
    We can combine terms and simplify if we change the range of the sum
    \begin{equation} \label{eq: 8.3}
        \Phi(r,\theta, \phi) = \sum_{l=0}^{\infty} \sum_{m=-l}^l (A_{l,m} r^l + B_{l,m} r^{-l-1}) e^{i m \phi} P_l^m (cos(\theta))
    \end{equation}            
    We would like to express the latter half of this equation using two orthonormal functions, one being a variable of $\theta$, indexed by l, and fixed by m, and the other being a variable of $\phi$ and indexed by m. This leads use to assign 
    \begin{equation} \label{eq: 8.4}
        Y_{l}^{m}(\theta, \phi) = a e^{i m \phi} P_l^m (cos(\theta))
    \end{equation}
    Where $a$ is an unknown normalization constant
    \newline Normalize to delta function
    \begin{equation} \label{eq: 8.5}
        I = \int_0^{2\pi} \int_0^{\pi} Y_{l}^{m} Y_{k}^{n*} sin(\theta) d\theta d\phi = \delta_{lk} \delta_{mn}
    \end{equation}
    Separate this integral over $\theta$ and $\phi$ because they are orthonormal
    \begin{equation} \label{eq: 8.6}
        \int_0^{2\pi} e^{i m \phi} e^{-i m \phi} d\phi = 2 \pi \delta_{mn}
    \end{equation}
    \begin{equation} \label{eq: 8.7}
        \int_0^{\pi} P_l^m (cos(\theta)) P_k^m cos(\theta)) sin(\theta) d\theta = \frac{2 (l+m)!}{(2l+1)(l-m)!}\delta_lk
    \end{equation}
    Combine results
    \begin{equation} \label{eq: 8.8}
        I = a^2 \frac{4 \pi (l+m)!}{(2l+l)(l-m)!} \delta_{lk} \delta_{mn}
    \end{equation}
    The appropriate value of a is one that makes $I=\delta_{lk} \delta_{mn}$; therefore they can be removed
    \newline Isolate a
    \begin{equation} \label{eq: 8.9}
        a = \sqrt{\frac{2l+l}{4\pi}} \sqrt{\frac{(l-m)!}{(l+m)!}}
    \end{equation}
    Substitute equation \ref{eq: 8.9} into equation \ref{eq: 8.4}
    \begin{equation} \label{eq: 8.10}
        Y_{l}^{m}(\theta, \phi) = \sqrt{\frac{2l+l}{4\pi}} \sqrt{\frac{(l-m)!}{(l+m)!}} e^{i m \phi} P_l^m (cos(\theta))
    \end{equation}
    Substitute the spherical harmonic variable from equation \ref{8.10} into the formulation from equation \ref{eq: 8.3}
    \begin{equation} \label{eq: 8.11}
        \Phi(r,\theta, \phi) = \sum_{l=0}^{\infty} \sum_{m=-l}^l (A_{l,m} r^l + B_{l,m} r^{-l-1}) Y_{l}^{m}(\theta, \phi)
    \end{equation}     
    Using equation \ref{eq: 8.11}, we can define the orthonormality of the spherical harmonics explicitly
    \begin{equation} \label{eq: 8.12}
        \int_0^{2\pi} \int_0^{\pi} Y_{l'm'}^{*}(\theta,\phi) Y_{lm}(\theta, \phi) sin(\theta) d\theta d\phi = \delta_{l'l} \delta_{m'm}
    \end{equation}
    Because the spherical harmonics form an orthonormal set of equations, and arbitrary function ($(\theta,\phi)$) can be expanded in terms of spherical harmonics
    \begin{equation} \label{eq: 8.13}
        f(\theta,\phi) = \sum_{l=1}^{\infty} \sum_{m=-l}^{l} A_{lm} Y_{lm} (\theta,\phi)
    \end{equation}
    Where 
    \begin{equation} \label{eq: 8.14}
        A_{lm} = \int_0^{2\pi} \int_0^{\pi} f(\theta,\phi) Y_{lm}^{*} (\theta, \phi) sin(\theta) d\theta d\phi
    \end{equation}

\subsection{Boundary Value and special case Spherical Harmonics}
    Using the Spherical Harmonic equation (Reference equation \ref{eq: 8.10}
    \begin{gather*}
        Y_{l}^{m}(\theta, \phi) = \sqrt{\frac{2l+l}{4\pi}} \sqrt{\frac{(l-m)!}{(l+m)!}} e^{i m \phi} P_l^m (cos(\theta))
    \end{gather*}
    For the case that $m=0$
    \begin{equation} \label{eq: 9.1}
        Y_{l}^{0}(\theta, \phi) = \sqrt{\frac{2l+l}{4\pi}} \sqrt{\frac{(l-0)!}{(l+0)!}} e^{i 0 \phi} P_l^0 (cos(\theta))
    \end{equation}
    \begin{equation} \label{eq: 9.2}
        Y_{l}^{0}(\theta, \phi) = \sqrt{\frac{2l+l}{4\pi}} P_l^0 (cos(\theta))
    \end{equation}
    For the case that $m=l$
    \begin{equation} \label{eq: 9.3}
        Y_{l}^{l}(\theta, \phi) = \sqrt{\frac{2l+l}{4\pi}} \sqrt{\frac{(l-l)!}{(l+l)!}} e^{i l \phi} P_l^l (cos(\theta))
    \end{equation}
    \begin{equation} \label{eq: 9.4}
        Y_{m}^{m}(\theta, \phi) = \frac{{(-1)}^m}{2^m m!} \sqrt{\frac{(2m+l)(2m)!}{4\pi}}  e^{i m \phi} (sin^m(\theta))
    \end{equation}
    For the case that $m<0$
    \begin{equation} \label{eq: 9.5}
        Y_{l}^{-m}(\theta, \phi) = {(-1)}^m Y_l^m*(\theta, \phi)
    \end{equation}
    For the case that $\theta=0$, if $m\neq0$
    \begin{equation} \label{eq: 9.6}
        Y_{l}^{m}(0, \phi) = 0
    \end{equation}
    For the case that $\theta=0$, if $m=0$
    \begin{equation} \label{eq: 9.7}
        Y_{l}^{m}(0, \phi) = \sqrt{\frac{2l + l}{4 \pi}}
    \end{equation}    
    For the case that $\theta = \pi$ and $m\neq0$
    \begin{equation} \label{eq: 9.8}
        Y_{l}^{m}(\pi, \phi) = 0
    \end{equation}
    For the case that $\theta=\pi$, if $m=0$
    \begin{equation} \label{eq: 9.9}
        Y_{l}^{m}(\pi, \phi) = (-1)^l \sqrt{\frac{2l + l}{4 \pi}}
    \end{equation}   
    For the case of a bounded sphere with a potential at the origin ($\Phi(r=r_0, \theta, \phi) = V(\theta, \phi)$
    \newline We begin with the electric potential for a spherical harmonic (Reference equation \ref{eq: 8.11}
    \begin{gather*}
        \Phi(r,\theta, \phi) = \sum_{l=0}^{\infty} \sum_{m=-l}^l (A_{l,m} r^l + B_{l,m} r^{-l-1}) Y_{l}^{m}(\theta, \phi)
    \end{gather*}
    Apply boundary condition
    \begin{equation} \label{eq: 9.10}
        V(\theta, \phi) = \sum_{l=0}^{\infty} \sum_{m=-l}^l (A_{l,m} r^l + B_{l,m} r^{-l-1}) Y_{l}^{m}(\theta, \phi)
    \end{equation}
    The region includes the origin, therefore, $B_{l,m}=0$ so that the solution does not explode to infinity
    \begin{equation} \label{eq: 9.11}
        V(\theta, \phi) = \sum_{l=0}^{\infty} \sum_{m=-l}^l (A_{l,m} r^l) Y_{l}^{m}(\theta, \phi)
    \end{equation}
    The solution is 
    \begin{equation} \label{eq: 9.12}
        \Phi(r,\theta, \phi) = \sum_{l=0}^{\infty} \sum_{m=-l}^l (A_{l,m} (\frac{r}{r_0})^l) Y_{l}^{m}(\theta, \phi)
    \end{equation}
    Where 
    \begin{equation} \label{eq: 9.13}
        A_{lm} = \int_0^{2\pi} \int_0^\pi V(\theta, \phi) Y_{lm}*(\theta,\phi) sin(\theta) d\theta d\phi
    \end{equation}

\end{document}










    
